\documentclass[a4paper,12pt]{article}
\usepackage[T2A]{fontenc}
\usepackage[utf8]{inputenc}
\usepackage[russian]{babel}
\usepackage{amsmath,amssymb,mathtools}
\renewcommand{\familydefault}{\sfdefault}
\usepackage{icomma}
\usepackage[a4paper,left=22mm,right=22mm,top=22mm,bottom=23mm]{geometry}
\usepackage{microtype}
\usepackage{tikz}
\usetikzlibrary{calc,arrows.meta,positioning,shapes.geometric}
\usepackage{tkz-euclide}
\usepackage{booktabs,tabularx}
\usepackage{enumitem}
\usepackage{fancyhdr}
\usepackage{setspace}
\usepackage{xcolor}

\definecolor{accent}{HTML}{008080}
\definecolor{darkaccent}{HTML}{004D4D}
\definecolor{textgray}{HTML}{333333}
\definecolor{softgray}{HTML}{68777A}
\color{textgray}
\onehalfspacing
\setlength{\parindent}{0pt}
\setlength{\parskip}{5pt}
\setlist[itemize]{leftmargin=6mm,itemsep=2pt,topsep=2pt}
\setlist[enumerate]{leftmargin=8mm,itemsep=8pt,topsep=4pt}
\renewcommand{\today}{24.10.2026}
\pagestyle{fancy}
\fancyhf{}
\fancyhead[L]{\footnotesize\color{softgray} Геометрия \;|\; 7 класс}
\fancyhead[R]{\footnotesize\color{softgray} Кирилл Зиновьев}
\fancyfoot[C]{\footnotesize\color{softgray}\thepage}
\renewcommand{\headrulewidth}{0.3pt}
\setlength{\headheight}{16pt}
\newcommand{\mysection}[1]{\vspace{3pt}\noindent{\large\bfseries\color{darkaccent}#1}\par\vspace{2pt}}
\newcommand{\pointitem}[1]{\item[\(\square\)] #1}
\newcommand{\ang}[1]{\angle #1}

\begin{document}
\thispagestyle{empty}
\begin{center}
\vspace*{33mm}
{\small\color{accent}\bfseries МАТЕМАТИКА \;\textbullet\; 7 КЛАСС}\par\vspace{15mm}
{\Huge\bfseries\color{darkaccent} Чек-лист}\par\vspace{9mm}
{\LARGE Равнобедренный треугольник\\[2mm]
и первый признак равенства треугольников}\par\vspace{20mm}
{\large Кирилл Зиновьев}\par\vspace{8mm}
{\normalsize \today}\par\vfill
{\small Лёвин Артём Александрович\\[2mm]
эксклюзивно для Кирилла Зиновьева}\par\vspace{11mm}
\begin{tikzpicture}[x=1cm,y=1cm]
\draw[accent,opacity=.65,line width=1pt]
(-6,0) .. controls (-3.3,.8) and (-2.1,-.6) .. (0,.12)
.. controls (2.1,.78) and (3.9,-.55) .. (6,0);
\end{tikzpicture}
\end{center}
\clearpage

\mysection{1. Что необходимо знать}
\begin{itemize}[label={}]
\pointitem{Равнобедренный треугольник имеет две равные стороны (боковые). Третья сторона называется основанием.}
\pointitem{Углы при основании равнобедренного треугольника равны.}
\pointitem{Первый признак равенства треугольников: если две стороны и угол \emph{между ними} одного треугольника соответственно равны двум сторонам и углу между ними другого, то треугольники равны.}
\pointitem{Общая сторона двух треугольников принадлежит каждому из них: например, в $\triangle BCD$ и $\triangle EDC$ сторона $CD$ общая.}
\pointitem{В равных треугольниках соответствующие стороны и углы равны. В записи равенства треугольников важен порядок вершин.}
\pointitem{Если две пары отрезков равны, то их разности равны: из $AC=AD$ и $AB=AE$ при $A-B-C$, $A-E-D$ следует $BC=DE$.}
\pointitem{Если соответствующие углы при пересечении двух прямых секущей равны, то эти прямые параллельны.}
\pointitem{Угол обозначают тремя буквами; его вершина стоит в середине: $\ang{CAD}$. Отрезок обозначают двумя буквами: $CD$.}
\end{itemize}

\mysection{2. Основная конфигурация}
Пусть $B$ лежит \emph{внутри} отрезка $AC$, а $E$ --- внутри $AD$:
\[
A-B-C,\qquad A-E-D,\qquad AC=AD,\qquad AB=AE.
\]
Слова «внутри отрезка» означают «строго между его концами».

\begin{center}
\begin{tikzpicture}[scale=.94]
\tkzDefPoint(0,3.75){A}
\tkzDefPoint(-3.25,0){C}
\tkzDefPoint(3.25,0){D}
\tkzDefPoint(-1.3,2.25){B}
\tkzDefPoint(1.3,2.25){E}
\tkzDrawSegments[line width=1pt](A,C A,D C,D)
\tkzDrawSegments[color=accent,line width=1.2pt](B,E)
\tkzDrawPoints[size=3](A,B,C,D,E)
\tkzLabelPoint[above](A){$A$}
\tkzLabelPoint[left](B){$B$}
\tkzLabelPoint[left](C){$C$}
\tkzLabelPoint[right](E){$E$}
\tkzLabelPoint[right](D){$D$}
\end{tikzpicture}
\end{center}
Здесь $\triangle ABE$ и $\triangle ACD$ равнобедренные, причём
\[
\ang{BAE}=\ang{CAD},\qquad
\ang{ABE}=\ang{AEB},\qquad
\ang{ACD}=\ang{CDA}.
\]

\mysection{3. Получение равных отрезков}
Поскольку точки располагаются в указанном порядке,
\[
BC=AC-AB,\qquad DE=AD-AE.
\]
Из $AC=AD$ и $AB=AE$ получаем $\boxed{BC=DE}$.

\textbf{Пример.} Если $AC=AD=18$ см и $AB=AE=7$ см, то
\[
BC=DE=18-7=11\text{ см}.
\]

\mysection{4. Как доказывать равенство треугольников}
\textbf{Пример 1.} Докажем, что $\triangle ABD\cong\triangle AEC$.

\begin{enumerate}[label=\arabic*.,itemsep=2pt]
\item $AB=AE$ и $AD=AC$ по условию.
\item $\ang{BAD}=\ang{EAC}$, так как это один и тот же угол с вершиной $A$.
\item Значит, $\triangle ABD\cong\triangle AEC$ по первому признаку.
\end{enumerate}
Соответствие вершин: $A\leftrightarrow A$, $B\leftrightarrow E$, $D\leftrightarrow C$.
Поэтому $\boxed{BD=EC}$ и $\ang{ABD}=\ang{AEC}$.

\textbf{Пример 2.} Докажем, что $\ang{CBD}=\ang{CED}$.

\begin{enumerate}[label=\arabic*.,itemsep=2pt]
\item Уже установлено, что $BC=DE$.
\item $CD=DC$ --- общая сторона треугольников $BCD$ и $EDC$.
\item $\ang{BCD}=\ang{EDC}$ --- углы при основании равнобедренного $\triangle ACD$.
\item Следовательно, $\triangle BCD\cong\triangle EDC$ по первому признаку.
\item Соответствующие углы при $B$ и $E$ равны:
$\boxed{\ang{CBD}=\ang{DEC}=\ang{CED}}$.
\end{enumerate}

\textbf{Важно.} В записи $\triangle BCD\cong\triangle EDC$ соответствуют вершины
$B\leftrightarrow E$, $C\leftrightarrow D$, $D\leftrightarrow C$.
Случайная перестановка вершин делает запись соответствия ошибочной.

\mysection{5. Как доказать параллельность $BE$ и $CD$}
Треугольники $ABE$ и $ACD$ равнобедренные, а их углы при вершине $A$ одинаковы. Поэтому их углы при основаниях равны:
\[
\ang{ABE}=\ang{ACD}.
\]
Поскольку $A$, $B$, $C$ лежат на одной прямой, это соответственные углы при секущей $AC$. Отсюда
\[
\boxed{BE\parallel CD}.
\]
\textbf{Не путайте:} отрезок $BE$ в общем случае \emph{не является средней линией} $\triangle ACD$: серединами сторон точки $B$ и $E$ могут и не быть.

\clearpage

\begin{center}
{\Large\bfseries\color{darkaccent} Алгоритм: доказательство равенства углов}\par\vspace{17mm}
\begin{tikzpicture}[
  >=Stealth,
  startstop/.style={ellipse,draw=darkaccent,line width=.8pt,align=center,text width=5.9cm,minimum height=1cm,fill=white,font=\small},
  act/.style={rectangle,rounded corners=2pt,draw=accent,line width=.85pt,align=center,text width=6.2cm,minimum height=1.25cm,fill=white,font=\small},
  test/.style={diamond,aspect=2.2,draw=darkaccent,line width=.9pt,align=center,text width=4cm,inner sep=2pt,fill=white,font=\small},
  helper/.style={rectangle,rounded corners=2pt,draw=softgray,line width=.85pt,align=center,text width=3.9cm,minimum height=2.3cm,fill=white,font=\small},
  link/.style={-Stealth,line width=.85pt,color=darkaccent}
]
\node[startstop] (start) {Нужно доказать\\равенство двух углов};
\node[act,below=13mm of start] (select) {Найдите треугольники, содержащие эти углы};
\node[test,below=15mm of select] (decide) {Есть две пары равных сторон и угол между ними?};
\node[helper,left=24mm of decide] (missing) {Найдите недостающие\\равенства:\\общую сторону,\\разности отрезков,\\углы при основании};
\node[act,below=17mm of decide] (sas) {Примените первый признак равенства треугольников};
\node[act,below=12mm of sas] (match) {Сопоставьте вершины и выберите нужные углы};
\node[startstop,below=12mm of match] (done) {Равенство углов доказано};
\draw[link] (start)--(select);
\draw[link] (select)--(decide);
\draw[link] (decide.south)--node[right,font=\small]{да}(sas.north);
\draw[link] (decide.west)--node[above,font=\small]{нет}(missing.east);
\draw[link] (missing.north) |- ($(select.west)+(-4mm,0)$) -- (select.west);
\draw[link] (sas)--(match);
\draw[link] (match)--(done);
\end{tikzpicture}
\end{center}
\clearpage

\mysection{Тренировочный блок --- путь А}
\textbf{10 заданий.} Выполняйте задания последовательно: от узнавания свойств к полноценным доказательствам.

\begin{enumerate}[label=\textbf{\arabic*.},itemsep=11pt]
\item В $\triangle PQR$ известно, что $PQ=PR$. Укажите основание и два равных угла.

\item В $\triangle ACD$ выполнено $AC=AD$, $\ang{CAD}=44^\circ$. Найдите $\ang{ACD}$ и $\ang{CDA}$.
\end{enumerate}

\textbf{Для заданий 3--10:} точки $B$ и $E$ лежат соответственно внутри отрезков $AC$ и $AD$; $AC=AD$, $AB=AE$. Пользуйтесь конфигурацией на странице~2.

\begin{enumerate}[label=\textbf{\arabic*.},start=3,itemsep=11pt]
\item Пусть $AC=AD=20$ см, $AB=AE=7$ см. Найдите $BC$ и $DE$.

\item Пусть $AC=AD=16$ см, $BC=DE=9$ см. Найдите $AB$ и $AE$.

\item Известно, что $\ang{CAD}=64^\circ$. Найдите $\ang{ABE}$ и $\ang{ACD}$. Обоснуйте равенство этих углов.

\item Рассмотрите $\triangle BCD$ и $\triangle EDC$. Укажите их общую сторону и пару равных углов, заключённых между сторонами $BC$, $CD$ и $DE$, $DC$ соответственно.

\item Запишите соответствие всех вершин для треугольников $ABD$ и $AEC$. Какое равенство третьих сторон следует из их равенства?

\item Докажите, что $BE\parallel CD$, используя свойства равнобедренных треугольников и признак параллельности прямых.

\item Докажите, что $\triangle ABD\cong\triangle AEC$ по первому признаку, и выведите равенство $\ang{ABD}=\ang{AEC}$.

\item Докажите, что $\ang{CBD}=\ang{CED}$, предварительно установив равенство двух подходящих треугольников по первому признаку.
\end{enumerate}
\clearpage

\mysection{Ответы и ориентиры для самопроверки}
\renewcommand{\arraystretch}{1.35}
\begin{tabularx}{\linewidth}{@{}>{\bfseries}p{9mm}X@{}}
\toprule
№ & Ответ или ключевой шаг доказательства \\
\midrule
1 & Основание $QR$; $\ang{PQR}=\ang{QRP}$.\\
2 & $\ang{ACD}=\ang{CDA}=68^\circ$.\\
3 & $BC=DE=13$ см.\\
4 & $AB=AE=7$ см.\\
5 & Оба угла по $58^\circ$: основания двух равнобедренных треугольников, угол при $A$ общий.\\
6 & Общая сторона $CD$; $\ang{BCD}=\ang{EDC}$.\\
7 & $A\leftrightarrow A$, $B\leftrightarrow E$, $D\leftrightarrow C$; $BD=EC$.\\
8 & $\ang{ABE}=\ang{ACD}$; это соответственные углы при секущей $AC$, поэтому $BE\parallel CD$.\\
9 & $AB=AE$, $AD=AC$, $\ang{BAD}=\ang{EAC}$; значит, $\triangle ABD\cong\triangle AEC$ и $\ang{ABD}=\ang{AEC}$.\\
10 & $BC=DE$, $CD$ --- общая, $\ang{BCD}=\ang{EDC}$; значит, $\triangle BCD\cong\triangle EDC$ и $\ang{CBD}=\ang{CED}$.\\
\bottomrule
\end{tabularx}
\end{document}
